\scanpage{190}
\mspage{64}
% archive sheet number: 189

s'égaliser par une flèche $\unsure{\varphi|\phi}(h)$, avec $h:i\longrightarrow j$ \deleted{\unsure{pour|dans}} Mais en fait \deleted{\illegible{2}} pour un $i'\in\mathrm{Ob}\,I'$ i.e un $Z\longrightarrow Y$ donné ($Z\in M^{\unsure{\pi|\Pi}})$ et un $i\in\mathrm{Ob}\,I$ i.e. un $Z_\alpha\lhook\joinrel\longrightarrow Y$ donné ($Z_\alpha\in\mathcal M^\pi$), il existe au plus \uline{une} flèche de \deleted{\unsure{$I'$}, $I=\unsure{M}$} dans $I'=\underline M^\pi/Y$ de $Z\longrightarrow Y$ dans $Z_\alpha\lhook\joinrel\longrightarrow Y$. Donc la condition F2 est trivialement satisfaite (on prend \deleted{\unsure{$j$}} $h=\mathrm{id}_i$).

\uline{Corollaire 5.11.} Si $(X_i)_{\unsure{i'|i}\in I}$ est un syst. inductif filtrant, avec les $X_i$ dans $M^\pi$, alors la catégorie $C=\underline M^\pi/M\times_{\underline{M}}(I,\varphi)$ construite p. 62 est filtrante, et le foncteur canonique \unsure{(fibrant)} $C\longrightarrow I$ est cofinal.

\deleted{Cela se} Que $C$ soit filtrante \inserted{et $C\to I$ cofinal,} résulte \unsure{immédiatement|également} du fait que $I$ l'est, \deleted{et} que les fibres de $C$ le sont, et que $C$ est cofibrée sur $I$. (Mérite un lemme, à \deleted{\unsure{toutes}} fins utiles.) \deleted{\unsure{Donc}}

\deleted{\unsure{On a maintenant}} \uline{NB} On peut \unsure{remplacer|compléter} $C$ par $C_0$ dont les fibres sont les $\mathrm{Ss}^{M^\pi}_{\pi}(X_i)\subset\underline M^{\pi}$; \unsure{à partir de|cofinale de} $C_i=\underline M^\pi/X_i$,

